% ======================================================================
\section{Disclosures and threats to validity}\label{sec:disclosures}

This section collects in one place everything a reader needs to weigh the numbers. Each item says what happened and
what it means for the claims.

\begin{enumerate}[label=\textbf{D\arabic*.}, leftmargin=*]
  \item \textbf{The test set is a seen set.} The test set (88 clips) was read after every accepted change -- by
        count about 19 times between 30 Sept and 2 Oct (one pre-registered read, one read after each of 13 accepted
        changes, and reads of five candidates that passed the development rule but were not accepted), plus re-reads of
        saved test pictures for sensitivity tables. The 60 older clips of the test set had already
        been read 11 times before the merge. No choice was made to improve a test number, with two exceptions that were
        undone (next item); still, choosing among many development variants and reading the test set after each makes
        the test numbers optimistic. They are a report on a seen set, not a blind confirmation.
  \item \textbf{Two test-informed choices, both undone.} A 4-s limit for the grouping rule was set after reading the
        test result of the unlimited version and was reverted to the pre-registered 8\,s; a picture-confidence floor
        suggested by a test read was not adopted. Rejections of three development passes (two ``expected sound'' rules and
        a second-listener rule) cited their test reads alongside screening-set and development reasons.
  \item \textbf{One-sided $p$-values.} All $p$-values of the final development phase are one-sided bootstrap shares (for
        example 0.034 one-sided is about 0.07 two-sided). The direct test of the final system against the September
        baseline has not been run (Section~\ref{sec:headline}).
  \item \textbf{Split overlap in September.} A second split file used by the stage-level detector bench
        shared 35 of its development clips with the 60 older test clips, and the FlexSED bar 0.8 was chosen on that
        split. On the clean development clips the 0.8 bar passes two of three adoption rules
        and misses the third by 0.011.
  \item \textbf{One annotator, two visibility re-checks.} Every label is the author's; no inter-rater agreement exists.
        The author re-checked visibility twice on the development set (3 of 14 disputed sounds changed, 0 of 20 random
        controls changed after a correction); the re-check targeted sounds where the pipeline and the gold disagreed,
        so it can favour the pipeline. Test labels were never changed.
  \item \textbf{AudioSet clips labelled with AudioSet labels visible.} The 19 AudioSet-Strong clips (across both sets)
        were labelled with the AudioSet labels available behind a button (shown on
        16 of 19), and DASM's checkpoint was selected on that AudioSet split. The labelling tool showed detector suggestions
        on 37 of the 50 newest clips, and older clips were pre-filled from the system under test.
  \item \textbf{``Wrong'' is strict.} A wrong picture includes redundant pictures of on-screen sounds and late second
        pictures of a long sound, not only false ones. A blind check found that about 1 in 6 ``wrong'' pictures is a real,
        wanted, off-screen sound missing from the gold. Absolute costs are pessimistic for every system.
  \item \textbf{The cost weights are post hoc.} $\beta = 2$ was declared for a gate sweep in September but became the
        reported metric only after per-sound F1 came out null; it was assumed, not measured with viewers.
  \item \textbf{Scored and live configurations differ in four known ways} (kinship direction, picture end cap, subject
        path, picture model; Section~\ref{sec:parity}). The kinship fix in particular was not re-scored.
  \item \textbf{Three development sounds are unreachable} by construction (labels the filter never draws), and the
        ceiling analysis was run on an earlier version of the final system.
  \item \textbf{The picture checker is not calibrated against a person.} Earlier panel rulings said no automatic
        picture checker should reject pictures until it matched the author's ratings; the final check-and-redraw loop
        was validated on known cases only and uses the same VLM that writes the subject. The planned sealed human
        confirmation of the pictures was cancelled. The picture claim rests on one blind glance test by the author.
  \item \textbf{The author is also the rater} of the glance test and the missing-sound checks. A second annotator packet
        exists but was not returned.
  \item \textbf{No study with deaf or hard-of-hearing viewers was run.} Every claim about benefit to viewers is
        modelled through the cost weights.
  \item \textbf{Larger vision model not tested.} Qwen3-VL-235B, the one lever every panel named for the gate, did
        not fit the available disk and is future work only.
  \item \textbf{Inconsistent records found while writing this report.} The pre-registration log is long (546\,KB) and
        was written under time pressure; the fact-check for this report found small inconsistencies, none of which
        changes an accepted change: two records of an intermediate version's test row (23/28 and 22/26, both cost 2.545,
        from before and after the grouping limit was reverted), two records of its development wrong split (6/9/3 and
        6/10/2, same total), a
        stale configuration comment (``$\leq$4\,s'' for an 8-s value), and the thesis chapter drafts of 27 Sept, which
        describe an earlier system and are superseded by this report.
\end{enumerate}
